% Polyhedral scheduling

\section{Scheduling SDGs with the Polyhedral Model}
\label{sec:analysis}

\begin{figure}[t]
  \centering
  \includegraphics[width=\columnwidth, trim={1 7 1 2}, clip]{cr_figures/sched_viz}
  %Use an example figure
  \caption{How \sys{} schedules $y[t] = f(x[t\!:\!min(t\!+\!3, T)])$}  
  \label{fig:scheduling_viz}
\end{figure}

RTs are declarative and thus require scheduling to find an execution order. Due to their dynamic dependencies, existing dataflow scheduling techniques~\cite{abadi2016tensorflow,moritz2018ray}, however, are unsuitable. \sys{} uses a \emph{polyhedral model}~\cite{bondhugula2008pluto,verdoolaege2010isl} to address this challenge. Scheduling occurs in two rounds: \sys first finds a schedule for tensor operations~(\S\ref{sec:polyhedral_model}), and then uses it to schedule memory management operations~(\S\ref{sec:mem_sched}).

\subsection{Polyhedral execution scheduling}
\label{sec:polyhedral_model}

%Scheduling a RL program is non-trivial due to their dynamic and cyclic dependencies.
%The polyhedral model is a powerful mathematical framework for program analysis that has its roots in scheduling dynamic uniform recurrence equations~\cite{karp1967organization}.
%While the polyhedral model has evolved to support a richer set of computations~\cite{benabderrahmane2010polyhedral,verdoolaege2010isl,bondhugula2008pluto}, today it is mainly adopted for optimizing small loop nests in imperative programs~\cite{grosser2012polly,vasilache2018tensor}, due to difficulty in modeling arbitrary programs.

The polyhedral model~\cite{verdoolaege2010isl, bondhugula2008pluto} has its origins in the scheduling of uniform recurrence equations~\cite{karp1967organization}, but has since evolved to support more complex computation~\cite{benabderrahmane2010polyhedral}. Today, it is used to optimize localized loop nests (kernels) in imperative programs~\cite{grosser2012polly, vasilache2018tensor}, as using it to model arbitrary programs remains challenging~\cite{verdoolaege2012polyhedral}. Since RTs are inspired by recurrence equations, \sys{} can apply the polyhedral model for whole-program scheduling. 
%\sys{} uses the \texttt{isl} library~\cite{verdoolaege2010isl} to schedule SDGs using the polyhedral model as shown in \F\ref{fig:scheduling_viz}.

%\noteps{Fairly sure we do not need this level of complexity. Try to reuse our own definitions of $\Omega$ and $\phi$}
%In a polyhedral model, each tensor~$x$ has its \emph{domain} represented by a polyhedron, $\Omega_{x} = \{ p \in \mathbb{Z}^{n} \; | \; \mathrm{constraints} \}$, which is the set of integer points in an $n$-dimensional space that satisfy a set of \emph{inequality constraints}.  A dependence of tensor~$x$ on $y$ is also represented by a polyhedron in the product of the two domains, $\phi_{x \rightarrow y} = \{(p,p') \in \Omega_{x} \times \Omega_{y} \; | \; \mathrm{constraints} \}$. For simplicity, we treat $\phi_{x \rightarrow y}$ as a function $\phi_{x \rightarrow y}(p) = \{p' \in \Omega_y | (p,p') \in \phi_{x \rightarrow y}\}$ that, given a point in $\Omega_x$, yields all points in $\Omega_y$ that it depends on.

%Additionally because of this, \sys can avoid complex dataflow analysis~\cite{grosser2012polly} to find hidden dependencies (e.g. caused by aliasing), which is typically the first step in polyhedral scheduling.

Scheduling in a polyhedral model finds a \emph{schedule function}~$\theta$ that maps each $n$-dimensional timestep of an operation to an \emph{execution time} in~$\mathbb{Z}^{s}$, $s \geq n$, and where the lexicographic order\footnote{Lexicographic order compares tuples left to right: $(1,2) \prec (1,5) \prec (2,1)$.} ($\prec$) determines what operation executes first. A \emph{valid} schedule must respect all dependencies, \ie  for $y[t]\!=\!f(x[\phi(t)])$ then $\forall t \in \Omega(y) . \theta(x[\phi(t)]) \prec \theta(y[t])$. 
Finding a schedule function requires solving an integer linear program~(ILP)~\cite{verdoolaege2010isl}, with an ILP objective which minimizes the sum of \emph{dependence distances} $\theta(y[t]) - \theta(x[\phi(t)])$ constrained by a set of \emph{validity constraints} derived from program dependencies. This yields schedules with low makespan.
%Through isl~\cite{verdoolaege2010isl}, \sys{} uses the pluto~\cite{bondhugula2008pluto} scheduler, whose objective function maximizes parallelism and locality by minimizing the \emph{dependence distance} $\theta(t) - \theta(\phi(t))$.

\F\ref{fig:scheduling_viz} shows the scheduling approach in \sys. SDGs already represent tensor domains, dependencies and conditions explicitly using symbolic expressions, which can be translated easily into polyhedral validity constraints (\ding{202} in \F\ref{fig:scheduling_viz}). In addition, \sys{} creates \emph{proximity constraints}, which guide the scheduler towards schedules with higher temporal locality and smaller tensor lifetimes. \sys{} uses the \emph{Pluto} scheduler~\cite{bondhugula2008pluto} through the \emph{isl} library~\cite{verdoolaege2010isl} (\ding{203}). 
In \F\ref{fig:scheduling_viz}, the schedule found delays $y$'s execution by 3 steps.

Since \sys fuses large regions of the SDG, the number of operations is reduced, making the ILP tractable even for large programs. 
\sys{} stores the resulting schedule function~$\theta$ (\ding{204}) for scheduling memory management operations, as discussed next. We describe how \sys scans the schedule to produce an imperative program AST (\ding{205}) in \S\ref{sec:ast_build}.


\subsection{Memory management}
\label{sec:mem_sched}

\begin{figure}[t]
  \centering
  \begin{subfigure}[t]{0.38\columnwidth}
    \centering
    \includegraphics[width=\columnwidth, trim={0 1 1 0}, clip]{cr_figures/dealloc_augment}
    \caption{Deallocation}\label{fig:mem_man:dealloc}
  \end{subfigure}%
  \hfill
  \begin{subfigure}[t]{0.58\columnwidth}
    \centering
    \includegraphics[width=\columnwidth, trim={0 1 1 0}, clip]{cr_figures/swap_augment}
    \caption{Swap}\label{fig:mem_man:swap}
  \end{subfigure}%
  \caption{SDG memory management augmentations}
  \label{fig:mem_man_augment}
\end{figure}


Memory management is a key concern in DL systems~\cite{peng2020capuchin,rhu2016vdnn,wang2018superneurons}. Using the polyhedral model, \sys{} automates memory management concerns: (i)~deallocations, (ii)~GPU/CPU memory swapping, and (iii)~buffer donations. 
%

The key idea is to \emph{augment} the SDG with memory management operations (and constraining dependencies), which enable \sys{} to reuse the polyhedral scheduler to find a memory management schedule.

\mypar{Scheduling deallocations} 
\F\ref{fig:mem_man:dealloc} shows how \sys{} augments the SDG with deallocations for the tensor produced by $P$ and consumed by $C$. Since $P$ produces a tensor for each timestep in its domain, $\Omega_P$, the deallocation $\mathrm{Dealloc}$ shares this same domain. \sys{} interprets the execution of $\mathrm{Dealloc}[t]$ as the time to deallocate the tensor $P[t]$.

Deallocation must always happen after production, which \sys{} models with an identity dependence, $\mathrm{Dealloc} \rightarrow P$. 
Then, to also ensure $P[t]$ is deallocated only after the consumer has finished using it, \sys{} adds a dependence from $\mathrm{Dealloc}$ to $C$, with the inverse dependence expression $C$ has on $P$. 
For example, if $\phi\!=\!t-3$, then $\mathrm{Dealloc}[t]$ must wait for $C[t+3]$ to run, before it can deallocate $P[t]$.

\mypar{Scheduling GPU/CPU memory swapping} 
\F\ref{fig:mem_man:swap} shows how \sys{} augments the SDG with evict (GPU-to-CPU) and load (CPU-to-GPU) memory movement operations. 
\sys{} inserts an initial evict operation $\mathrm{Evict_A}$, which evicts the tensor $P[t]$ immediately after production.
It then inserts $\mathrm{Load}$ and $\mathrm{Evict_B}$ operators, to re-load and again evict all timesteps of $P$ needed by $C$ at timestep $t$. All operations have domain $\Omega_P$ and \sys{} interprets the execution of $\mathrm{Evict_{A/B}}[t]$ and $\mathrm{Load}[t]$ as the time to evict or load $P[t]$, respectively.


A tensor can only be evicted after it has been produced, and loaded after it has been evicted. \sys{} models this with the identity dependencies, $\mathrm{Load} \rightarrow \mathrm{Evict_A} \rightarrow P$.
To guarantee that the timesteps of $P$ needed by $C[t]$ are available in-memory when $C[t]$ executes, \sys{} adds $C \xrightarrow{\smash{\phi}} \mathrm{Load}$.
Finally, similarly to deallocations, $\mathrm{Evict_B}$ must execute only after the last use of $P[t]$ by $C$, so \sys{} adds the inverse of the consumer's dependence from $\mathrm{Evict_B}$ to $C$.


By placing these dependencies, \sys{} introduces constraints on when the memory management operations can execute, which guarantees that the resulting schedule is valid and efficient.
To prevent thrashing~\cite{denning1968working}, \sys{} does not currently swap frequently accessed tensors, such as DNN parameters or K/V caches in attention models.
At runtime, \sys{} performs such swap operations \emph{asynchronously}, through a shared \emph{page-locked} host-memory manager, which pre-allocates page-locked buffers for each swapped RT.

\mypar{Finding buffer donations} Given the SDG's stateless nature, it is also important to detect when a tensor produced by one operation can be \emph{donated} to be reused by another operation. This avoids expensive re-allocations and enables in-place mutation of, \eg DNN parameters.

After a tensor is donated, it is invalidated and can no longer be used by other operations.
Thus, an operator $O_d$ can only donate its output tensor to one of its dependents, called the \emph{donation receiver} $O_r$, if for every timestep $t$ in its domain $\Omega(O_d)$, $O_r$ is the last user of $O_d[t]$. This requires that it is scheduled \emph{after} all other competing dependents $O_c$ according to the schedule function~$\theta$. Formally, \sys donates $O_d$ to $O_r$ if and only if
\begin{align*}
\forall t \in \Omega(O_d) .\ \forall O_c \in \mathcal{D}(O_d) \setminus \{O_r\} .\ \theta\left(\phi^{-1}_{c}(t)\right) \prec \theta\left(\phi^{-1}_{r}(t)\right)
\end{align*}
where $\mathcal{D}(O_d)$ is the set of dependents of $O_d$, $\phi^{-1}_{c}$ and $\phi^{-1}_{r}$ are the inverse dependence expressions from $O_c$ and $O_r$ to $O_d$ respectively, and $\setminus$ is the set difference operator. %This information is used by \sys{}'s runtime when it is code-generating kernels (\S\ref{sec:runtime}), so that code-generators are aware that, for example, more efficient in-place operations can be used.

\subsection{Generating imperative programs}
\label{sec:ast_build}

\sys{} must obtain an imperative program for the user's declarative RT computation, but the schedule function~$\theta$ produced by the polyhedral scheduler~\cite{bondhugula2008pluto} is not directly executable. \sys{} transforms it into an abstract syntax tree~(AST) by using the polyhedral library~\cite{verdoolaege2010isl} to scan~\cite{grosser2015polyhedral} it. The AST contains the following types of statements:
%
\begin{myitemize}
  \item \textbf{sequence}: a block of AST nodes executed in order;%or in parallel.
  \item \textbf{if}: a conditional with then and else sub-ASTs;
  \item \textbf{loop}: a generic loop with an iteration counter, start value, condition, step and body sub-AST;
  %TODO: removed parallel due to space
  \item \textbf{execute}: executes operation~$O[t]$, where $t$ is given by a mapping from loop counters; and
  \item \textbf{deallocate/load/evict}: executes the memory instruction for timestep~$t$ (or multiple steps at once) given by a mapping from loop counters.
\end{myitemize}

To generate efficient ASTs, \sys applies further optimizations: (1)~it \emph{partially unrolls} loops to reduce the number of nested conditionals; (2)~it \emph{promotes} static loops that contain only memory management operations into batched instructions by mapping the loop's range to a slice of points; and (3)~it finds and removes redundant loads/evicts by removing (i)~any evict followed by a load or deallocation for the same tensor; and (ii)~loads for already in-memory tensors. %(as long as there is no user in between).


By interpreting the AST at runtime, users are able to step through the program with high granularity. This allows for simpler debugging and integration with existing DL libraries.%, but it can also be a performance bottleneck due to interpretation overheads for small programs, which we leave for future work.

% End
